% !TeX program = xelatex
\documentclass[UTF8,zihao=-4,fontset=fandol]{ctexart}
\newcommand{\ExamNumber}{2}
\usepackage{linear-algebra-exam}

\begin{document}
\ExamCover{7}{64}{50}

\Question{1}{单项选择}{2009 年数学一第 6 题、数学二第 7 题（均为原卷第 2 页）}{4}{***}
设 $A,B$ 均为 $2$ 阶矩阵，$A^*,B^*$ 分别为 $A,B$ 的伴随矩阵，若 $|A|=2,\ |B|=3$，则分块矩阵 $\begin{pmatrix}O&A\\ B&O\end{pmatrix}$ 的伴随矩阵为（\quad）。
\Choices
  {$\begin{pmatrix}O&3B^*\\2A^*&O\end{pmatrix}$}
  {$\begin{pmatrix}O&2B^*\\3A^*&O\end{pmatrix}$}
  {$\begin{pmatrix}O&3A^*\\2B^*&O\end{pmatrix}$}
  {$\begin{pmatrix}O&2A^*\\3B^*&O\end{pmatrix}$}

\Solution
\Known $O$ 表示 $2$ 阶零矩阵。记题中的 $4$ 阶矩阵为 $M$，要求它的伴随矩阵 $M^*$。伴随矩阵由各元素的代数余子式转置排列而成。
\Method 如果 $M$ 可逆，伴随矩阵可以由 $M^*=|M|M^{-1}$ 得到。题目给出了 $A,B$ 的非零行列式，因此可直接用 $A^{-1},B^{-1}$ 来构造分块矩阵 $M^{-1}$。
\Derive 将 $M$ 的第 $1$、$3$ 行互换，再将第 $2$、$4$ 行互换，得到分块对角矩阵 $\begin{pmatrix}B&O\\O&A\end{pmatrix}$。两次互换使行列式变号两次，所以
\[
|M|=|B|\,|A|=6.
\]
按分块乘法，$M\begin{pmatrix}O&B^{-1}\\A^{-1}&O\end{pmatrix}$ 的左右对角块分别是 $AA^{-1}$ 和 $BB^{-1}$，其余两块都是零。因此右边的分块矩阵就是 $M^{-1}$。再由 $B^{-1}=\frac{B^*}{3}$、$A^{-1}=\frac{A^*}{2}$，得到
\[
M^*=6M^{-1}=\begin{pmatrix}O&2B^*\\3A^*&O\end{pmatrix}.
\]
写逆矩阵时可从乘法结果倒推：分块矩阵 $M$ 左乘另一个矩阵，第一块行由 $A$ 乘该矩阵的第二块行，第二块行由 $B$ 乘第一块行。因此要得到 $E_4$，候选逆矩阵的左下块须为 $A^{-1}$，右上块须为 $B^{-1}$，两个对角块须为零。这里右上块乘 $6$ 再除以 $|B|=3$，得到的系数是 $2$；左下块则乘 $6$ 再除以 $|A|=2$，系数是 $3$。
\Verify 选 \textbf{B}。把结果左乘 $M$，两个对角块分别是 $3AA^*=6E$ 和 $2BB^*=6E$，得到 $|M|E_4$，其中 $E$、$E_4$ 分别为二阶、四阶单位矩阵。

\Question{2}{填空}{2021 年数学二第 16 题（原卷第 2 页）}{8}{****}
多项式
\[
f(x)=\begin{vmatrix}
x & x & 1 & 2x\\
1 & x & 2 & -1\\
2 & 1 & x & 1\\
2 & -1 & 1 & x
\end{vmatrix}
\]
中 $x^3$ 项的系数为 \blank。

\Solution
\Known $f(x)$ 是一个四阶行列式展开后的多项式。只需找 $x^3$ 项的系数。
\Method 根据行列式的定义，一个乘积项必须在每行、每列各取一个元素。要形成 $x^3$，四个因子中必须恰好三个含 $x$，另一个是常数。
\Derive 如果第一行取第 $1$ 列的 $x$，后面三行的含 $x$ 项都在主对角线上。一旦从这三行取了两个 $x$，最后所剩的位置还是对角线上的 $x$，所得次数为四。第一行若取第 $3$ 列的 $1$，它占去了第三行 $x$ 的列，无法从其余三行都取到 $x$。

剩下两种取法。第一行取第 $2$ 列的 $x$ 时，第二行只能取常数，第三、四行必须取各自对角线上的 $x$。列序为 $(2,1,3,4)$，逆序数为 $1$，贡献 $-x\cdot 1\cdot x\cdot x=-x^3$。

第一行取第 $4$ 列的 $2x$ 时，第四行只能取常数，第二、三行必须取各自对角线上的 $x$。列序为 $(4,2,3,1)$，逆序数为 $5$，贡献 $-(2x)\cdot x\cdot x\cdot2=-4x^3$。
可把符号计算得更明确。$(2,1,3,4)$ 只有 $2$ 在 $1$ 前面这一处逆序。$(4,2,3,1)$ 中，$4$ 与后面三个数各形成一个逆序，$2$、$3$ 又分别与最后的 $1$ 形成逆序，共五处。两种取法的常数因子分别是第 $(2,1)$ 元的 $1$ 和第 $(4,1)$ 元的 $2$。因此除去 $x^3$ 后，带符号的乘积是 $-1$ 与 $-(2\cdot2)=-4$。
\Verify 填 $\boxed{-5}$。两项分别选了第一行的两个非对角含 $x$ 元素，列序均为奇排列，符号均为负。

\Question{3}{填空}{2010 年数学二第 14 题（原卷第 2 页）}{8}{***}
设 $A,B$ 为 $3$ 阶矩阵，且 $|A|=3,\ |B|=2,\ |A^{-1}+B|=2$，则 $|A+B^{-1}|=\blank$。

\Solution
\Known 已知 $A,B$ 的行列式以及 $A^{-1}+B$ 的行列式，要求 $A+B^{-1}$ 的行列式。因 $|A|$、$|B|$ 均非零，两个逆矩阵都存在。
\Method 已知的行列式含 $A^{-1}$ 和 $B$，所求的含 $A$ 和 $B^{-1}$。试从左侧乘以 $A$、右侧乘以 $B^{-1}$，用互逆关系把中间的两项化为所求的两项。
\Derive 按矩阵乘法的结合律、分配律，
\[
A(A^{-1}+B)B^{-1}
=AA^{-1}B^{-1}+ABB^{-1}
=B^{-1}+A\text{。}
\]
因此可用乘积的行列式公式。又因 $BB^{-1}=E$，有 $|B^{-1}|=\frac{1}{|B|}=\frac{1}{2}$，从而
\[
|A+B^{-1}|=|A|\,|A^{-1}+B|\,|B^{-1}|
=3\cdot2\cdot\frac{1}{2}=3\text{。}
\]
也可通过一个公共因子检查选边是否正确。$A^{-1}+B=A^{-1}(E+AB)$，而 $A+B^{-1}=(E+AB)B^{-1}$。因为 $|A^{-1}|=\frac{1}{3}$，由第一式得 $|E+AB|=6$；第二式再乘 $|B^{-1}|=\frac{1}{2}$。两种分解使用同一个 $AB$，无须假设 $AB=BA$。
\Verify 填 $\boxed{3}$。把等式反向写成 $A^{-1}(A+B^{-1})B=A^{-1}+B$，取行列式后也得 $\frac{1}{3}\cdot3\cdot2=2$，符合已知条件。

\Question{4}{填空}{2013 年数学二第 14 题、数学一第 13 题（均为原卷第 2 页）}{8}{****}
设 $A=(a_{ij})$ 是 $3$ 阶非零矩阵，$|A|$ 为 $A$ 的行列式，$A_{ij}$ 为 $a_{ij}$ 的代数余子式。若 $a_{ij}+A_{ij}=0$（$i,j=1,2,3$），则 $|A|=\blank$。

\Solution
\Known 各元素与其代数余子式互为相反数。要求 $3$ 阶实矩阵 $A$ 的行列式，题设还保证 $A$ 至少有一个非零元素。
\Method 题设逐个约束余子式，先将它写成伴随矩阵恒等式。伴随矩阵的第 $(i,j)$ 个元素为 $A_{ji}$，所以用到的转置不能省略。
\Derive 由 $A_{ij}=-a_{ij}$ 得 $A^*=-A^{\mathrm T}$。设 $d=|A|$，利用 $AA^*=dE$ 可得
\[
AA^{\mathrm T}=-dE.
\]
左边的第 $i$ 个对角元素是第 $i$ 行各分量的平方和。至少有一行非零，因此至少有一个这样的平方和为正，从而 $-d>0$。这一步保证接下来可以除以 $d^2$。

取等式两边的行列式。转置不改变行列式，右侧三阶对角矩阵的行列式为 $(-d)^3$，所以 $d^2=(-d)^3$。除以 $d^2$ 即得 $d=-1$。
从矩阵元素来看，$AA^{\mathrm T}$ 的第 $i$ 个对角元为 $a_{i1}^2+a_{i2}^2+a_{i3}^2$。在实数范围内，只要该行有非零元素，这个和就严格为正。若省略非零条件，$A=O$ 会满足题给元素关系且行列式为零。这里先用题设排除 $d=0$，再由 $d^2=-d^3$ 得 $1=-d$。
\Verify 填 $\boxed{-1}$。取 $A=\operatorname{diag}(-1,1,1)$，它的代数余子式矩阵确实为 $-A$，且行列式为 $-1$。

\Question{5}{填空}{2023 年数学一第 15 题（原卷第 2 页）}{8}{***}
已知向量
\[
\alpha_1=\begin{pmatrix}1\\0\\1\\1\end{pmatrix},\quad
\alpha_2=\begin{pmatrix}-1\\-1\\0\\1\end{pmatrix},\quad
\alpha_3=\begin{pmatrix}0\\1\\-1\\1\end{pmatrix},\quad
\beta=\begin{pmatrix}1\\1\\1\\-1\end{pmatrix},
\]
$\gamma=k_1\alpha_1+k_2\alpha_2+k_3\alpha_3$。若 $\gamma^{\mathrm T}\alpha_i=\beta^{\mathrm T}\alpha_i\ (i=1,2,3)$，则 $k_1^2+k_2^2+k_3^2=\blank$。

\Solution
\Known 三个列向量加权相加得到 $\gamma$。已知 $\gamma^{\mathrm T}$ 左乘这三列的结果分别等于 $\beta^{\mathrm T}$ 的相应乘积，要求权重的平方和。
\Method 把三列并成矩阵 $S$，将权重记为列矩阵 $k$。由 $\gamma=Sk$，可把题给的三个标量等式合写为
\[
k^{\mathrm T}S^{\mathrm T}S=\beta^{\mathrm T}S.
\]
\Derive 利用题给元素直接做行乘列，可得
\[
S^{\mathrm T}S=3E,\qquad \beta^{\mathrm T}S=(1,-3,-1).
\]
例如，$(1,2)$ 元是 $1\cdot(-1)+0\cdot(-1)+1\cdot0+1\cdot1=0$，而 $(1,1)$ 元是 $1^2+0^2+1^2+1^2=3$。因此
\[
k^{\mathrm T}=\left(\frac{1}{3},-1,-\frac{1}{3}\right).
\]
剩余的两个交叉乘积也可逐项检查。$\alpha_1^{\mathrm T}\alpha_3=1\cdot0+0\cdot1+1\cdot(-1)+1\cdot1=0$，$\alpha_2^{\mathrm T}\alpha_3=(-1)\cdot0+(-1)\cdot1+0\cdot(-1)+1\cdot1=0$。三个列自身的平方和均为 $3$，所以三个方程分别只剩 $3k_i$。最终的平方和为 $\frac{1}{9}+1+\frac{1}{9}=\frac{11}{9}$。
\Verify 填 $\boxed{\frac{11}{9}}$。把这个 $k$ 代回，$3k^{\mathrm T}=(1,-3,-1)$，同时满足题目给出的三个等式。

\Question{6}{填空}{2025 年数学一第 15 题（原卷第 2 页）}{8}{***}
设矩阵 $A=\begin{pmatrix}4&2&-3\\a&3&-4\\b&5&-7\end{pmatrix}$，若方程组 $A^2x=0$ 与 $Ax=0$ 不同解，则 $a-b=\blank$。

\Solution
\Known $Ax=0$ 的每个解都会使 $A^2x=0$。题设指出两者的解不全相同，要求参数之差。
\Method 如果 $A$ 可逆，从 $A^2x=0$ 左乘 $A^{-1}$ 就得到 $Ax=0$，两者会同解。这与题设矛盾，所以可以用 $|A|=0$ 求参数关系。
\Derive 按第一列展开，
\[
|A|=4(-21+20)-a(-14+15)+b(-8+9)=-4-a+b.
\]
令其为零，得到 $a-b=-4$。
第一列位置的符号为 $+,-,+$。删去第 $1$、$2$、$3$ 行得到的二阶余子式分别为 $-1,1,1$，所以三个带符号项是 $-4,-a,+b$。只需要证明原条件推出 $|A|=0$，再求差值。$|A|=0$ 本身不足以反推两个方程不同解，因为非零的 $Ax$ 还必须能被 $A$ 再次乘成零。
\Verify 填 $\boxed{-4}$。例如取 $a=5,\ b=9$，令 $x=(1,-1,0)^{\mathrm T}$，则 $Ax=(2,2,4)^{\mathrm T}\ne0$，再左乘 $A$ 得零，确实可以出现题设情况。

\Question{7}{解答}{2003 年数学一（见 2019 年数学一配套解析 PDF 第 9 页）}{20}{*****}
已知平面上三条不同直线的方程分别为
\[
l_1:\ ax+2by+3c=0,\qquad
l_2:\ bx+2cy+3a=0,\qquad
l_3:\ cx+2ay+3b=0.
\]
试证这三条直线交于一点的充分必要条件为 $a+b+c=0$。

\Solution
\Known 三条直线互不相同。交于一点意味着三式有同一组 $(x,y)$，并且这一组解唯一。需要证明交点条件与 $a+b+c=0$ 可以互推。
\Method 证明必要性时，有共同解意味着三式的常数列可由系数两列加权相加得到，所以三列组成的行列式为零。证明充分性时，利用三式的和为零，把问题化为前两式是否有唯一解。
\Derive 若有共同解，则在下列行列式中把第一列的 $x$ 倍、第二列的 $y$ 倍加到第三列，第三列变为零。因此
\[
0=\begin{vmatrix}a&2b&3c\\b&2c&3a\\c&2a&3b\end{vmatrix}
=-6(a+b+c)(a^2+b^2+c^2-ab-bc-ca).
\]
最后一个因式为 $\frac{1}{2}\big((a-b)^2+(b-c)^2+(c-a)^2\big)$。若它为零，则 $a=b=c$，三式相同，违背题设。所以 $a+b+c=0$。

反过来，若 $a+b+c=0$，三式左端相加恒为零。前两式的系数行列式为 $2(ac-b^2)=-2(a^2+ab+b^2)$。由于 $a^2+ab+b^2=(a+\frac{b}{2})^2+\frac{3b^2}{4}$，其为零只可能是 $a=b=0$，继而 $c=0$，此时无法给出三条不同直线。因此该系数行列式非零，前两式有唯一公共解，它也自动满足第三式。
三阶行列式的因式可以逐步核对。先分别从第二、三列提出 $2,3$，再按第一行展开，括号内为
\[
a(bc-a^2)-b(b^2-ac)+c(ab-c^2)=3abc-a^3-b^3-c^3.
\]
用 $a^3+b^3+c^3-3abc=(a+b+c)(a^2+b^2+c^2-ab-bc-ca)$ 得到原式。充分性一边，前两式左端加第三式左端恒为零，所以只要前两式成立，第三式就自动成立。按克拉默法则，唯一交点的两个坐标可以写成
\[
x=\frac{3(ab-c^2)}{ac-b^2},\qquad
y=\frac{3(bc-a^2)}{2(ac-b^2)}.
\]
\Verify 两个方向均成立。例如 $a=1,b=0,c=-1$ 时，前两式给出 $x=3,\ y=\frac{3}{2}$，代入第三式也成立。
\end{document}
